\documentclass[11pt]{article}

\usepackage[utf8]{inputenc}
\usepackage[T1]{fontenc}
\usepackage{lmodern}
\usepackage{geometry}
\usepackage{amsmath,amssymb,amsfonts,amsthm,mathtools}
\usepackage{bm}
\usepackage{enumitem}
\usepackage{microtype}
\usepackage{hyperref}
\usepackage{titlesec}
\usepackage{booktabs}
\usepackage{array}
\usepackage{setspace}
\usepackage{mathrsfs}
\usepackage{physics}

\geometry{margin=1in}
\hypersetup{colorlinks=true, linkcolor=black, urlcolor=blue, citecolor=black}
\setlist[enumerate]{topsep=0.4em,itemsep=0.35em}
\setlist[itemize]{topsep=0.3em,itemsep=0.25em}
\titleformat{\section}{\large\bfseries}{\thesection.}{0.5em}{}
\titleformat{\subsection}{\normalsize\bfseries}{\thesubsection.}{0.5em}{}
\setstretch{1.06}

\newcommand{\Aether}{\AE{}ther}
\newcommand{\Aflow}{\AE{}ther-flow}
\newcommand{\TheoryName}{\Aflow{} Interpretation of Relativity}
\newcommand{\TheoryProgram}{\Aether{} / \Aflow{} framework}
\newcommand{\CompletedMetric}{g^{\sharp}}

\newcommand{\ReducedStructure}{\mathfrak S^{\mathrm{disp,resp,red}}}
\newcommand{\CurrentBodyImageCore}{\mathfrak S^{\mathrm{disp,resp,cbimg}}}

\newtheorem{definition}{Definition}
\newtheorem{proposition}{Proposition}
\newtheorem{remark}{Remark}
\newtheorem{corollary}{Corollary}
\newtheorem{theorem}{Theorem}

\title{\textbf{The \TheoryName{}}\\[0.4em]
\large Primitive-Reservoir Observer-Reduction\\
\large Section-Centered Hidden-Response\\
\large Canonical Current-Body Image Core\\
\large Necessity / Uniqueness from Reduced Projection-Response\\
\large Substrate Structure Next Benchmark Criterion}
\author{}
\date{}

\begin{document}

\maketitle

\begin{abstract}
The current active-primary theorem now proves that, once the fixed current displacement body, the finite-dimensional substrate carrier, the surjective displacement projection, the coercive positive self-adjoint substrate-response operator, exact-shell discipline, and benchmark faithfulness are held fixed, the full canonical current-body image core \(\CurrentBodyImageCore\) is canonically derived and uniquely forced. The canonical image body, the restricted projection / canonical-lift relation, the canonical image response law, and the canonical exact-shell image are therefore no longer independently live benchmark-facing burdens inside the current image-core frontier.

This note records the routing consequence of that gain. The next honest benchmark-facing move must now derive or justify the reduced projection-response substrate structure \(\ReducedStructure\) itself from still more primitive explicit substrate structure, or prove a sharper theorem that makes that derivation unavoidable, while preserving the same benchmark one-metric package, the same ordinary matter route, and the same claim boundary. The one operative metric remains exactly \(\CompletedMetric\), there is no second operative metric, the low-energy matter law is not enlarged, and no stronger promotion language is licensed.
\end{abstract}

\section{Introduction}

The active derivational program of \TheoryName{} is no longer at body-level, relation-level, response-level, shell-level, or even full-image-core scope inside the current section-centered hidden-response frontier. The new full-image-core forcing theorem has already spent that move. Once the reduced projection-response substrate structure is fixed, the entire canonical current-body image core \(\CurrentBodyImageCore\) follows canonically and uniquely.

The present note records the routing consequence of reading that theorem together with the benchmark safeguard against recursive depth drift. It does not reopen the already-screened larger substrate-body presentation, the older projection-operator core, the full displacement-response substrate law, or the already-spent body-, relation-, response-, shell-, or full-image-core presentations as if any of them were again independently live. It records only which kind of next manuscript would now count as an honest benchmark-facing gain below the new frontier.

\paragraph{Framework and claim boundary.}
\TheoryProgram{} denotes the broader \Aether{} / \Aflow{} framework built on the claims that the \Aether{} is the underlying four-dimensional substrate of reality, the \Aflow{} is its intrinsic ordered motion, observed three-dimensional space is the local experiential slice of that deeper substrate, S-time is the experienced order of change arising from matter, light, and the \Aflow{}, and observed expansion is the three-dimensional appearance of deeper four-dimensional ordered motion. Within that framework, \TheoryName{} denotes the exact-closure theory in which the effective gravitational dynamics are adopted to be exactly Einsteinian with universal matter coupling. In the active sequence, \emph{adoption} means use of that established relativistic dynamics without claiming substrate derivation, while \emph{derivation} is reserved for a first-principles recovery from explicit substrate variables.

The present note stays strictly inside that boundary. It records only which kind of next manuscript would now count as an honest benchmark-facing gain below the canonical-current-body-image-core necessity / uniqueness frontier.

\section{Fixed Inputs}

\begin{definition}[Full-image-core forcing frontier inputs]
The criterion recorded here uses the following fixed inputs.
\begin{enumerate}
    \item The benchmark exact-closure package, which fixes the public one-metric GR target of the repository.
    \item The benchmark gatekeeping layer, including the recorded safeguard that blocks same-output deeper-origin relays from counting as new front-line progress once the live benchmark-relevant datum is unchanged.
    \item The preserved theorem chain showing that the canonical image body, the restricted projection / canonical-lift relation, and the canonical image response / exact-shell data were each separately forced inside the current image-core frontier.
    \item The current theorem proving that, once the reduced projection-response substrate structure \(\ReducedStructure\) is fixed, the full canonical current-body image core \(\CurrentBodyImageCore\) is canonically derived and unique at benchmark scope.
\end{enumerate}
\end{definition}

\begin{remark}
The present note does not replace those manuscripts. It records the routing consequence of reading them together under the active safeguard against recursive depth drift.
\end{remark}

\section{Why the Live Burden Now Sits Below Reduced Structure}

\begin{proposition}[The live burden is renewed below reduced-structure restatement]
At the current stage of the primitive-reservoir observer-reduction line, the full-image-core forcing theorem renews the live benchmark-facing burden below any mere restatement of the same reduced projection-response substrate structure \(\ReducedStructure\) or of the same canonical current-body image core \(\CurrentBodyImageCore\).
\end{proposition}

\begin{proof}
That theorem changes benchmark-facing status because it proves that, once the reduced projection-response substrate structure is fixed, the full image core itself is already forced. Once that derivation and uniqueness statement is in hand, another manuscript that merely preserves the same reduced structure and therefore the same image body, the same restricted projection / canonical-lift relation, the same image response law, the same exact-shell image, and the same downstream benchmark-facing consequences no longer removes any remaining benchmark-facing freedom. The live burden therefore no longer sits on restating the same full image core or merely renaming the same reduced structure. It sits on deriving that reduced structure from still more primitive explicit substrate structure, or on proving a sharper theorem that makes such a derivation unavoidable.
\end{proof}

\begin{definition}[Benchmark-neutral same-structure relay below the current frontier]
A \emph{benchmark-neutral same-structure relay below the current frontier} is a manuscript that preserves the same benchmark one-metric output, the same ordinary matter route, the same reduced projection-response substrate structure \(\ReducedStructure\), the same full canonical current-body image core \(\CurrentBodyImageCore\), and the same downstream benchmark-facing consequences while merely relocating the unexplained substrate layer deeper, renaming the same reduced structure, restoring the already-screened larger substrate-body presentation, or re-expanding back to the older projection-operator core, the full displacement-response substrate law, or the larger pullback-body-operator, pullback-operator, pullback-quadratic, hidden-response, residual-normal-form, or pre-proto-section presentations without a new reduced-structure derivation, new forcing theorem, or comparably sharp new gain.
\end{definition}

\section{The Current Post-Reduced-Structure Criterion}

\begin{definition}[Admissible next benchmark-facing manuscript below the current frontier]
Below the current canonical-current-body-image-core necessity / uniqueness frontier, a manuscript counts as an admissible next benchmark-facing move only if it does at least one of the following:
\begin{enumerate}
    \item derives or justifies the reduced projection-response substrate structure \(\ReducedStructure\) itself from still more primitive explicit substrate structure;
    \item proves a sharper necessity, uniqueness, or compatibility theorem that makes such a reduced-structure derivation unavoidable while preserving the same benchmark one-metric package, the same ordinary matter route, and the same claim boundary;
    \item does one of the previous two without reopening the larger substrate-body presentation, the older projection-operator core, the full displacement-response substrate law, or the already-spent full image-core package as if those were again the live burden.
\end{enumerate}
\end{definition}

\begin{theorem}[Next honest benchmark-facing task below the current frontier]
The next honest benchmark-facing manuscript below the current canonical-current-body-image-core necessity / uniqueness frontier must derive or justify the reduced projection-response substrate structure \(\ReducedStructure\) from still more primitive explicit substrate structure, or prove a still sharper theorem that makes such a derivation unavoidable. A manuscript that only repeats the same reduced structure or the same full image core through another deeper relay, restores the larger substrate-body presentation, re-expands the discussion back to the older projection-operator core or the full displacement-response substrate law, or replays the larger already-spent presentations without such a new gain does not satisfy the criterion.
\end{theorem}

\begin{proof}
By the preceding proposition, the live burden already lies below any mere restatement of the same reduced structure or the same full image core. Therefore a subsequent manuscript must improve on that now-forced package rather than merely accompany it. A deeper-origin manuscript can do so only if it changes benchmark-facing status by more than depth alone. If it simply reproduces the same reduced structure, the same full image core, and the same downstream outputs, the routing safeguard classifies it as benchmark neutral. The remaining honest alternatives are therefore exactly those stated in the definition.
\end{proof}

\section{What the Next Move Should Not Be}

\begin{proposition}[Screened next moves below the current frontier]
Below the current canonical-current-body-image-core necessity / uniqueness frontier, none of the following is the default next benchmark-facing move:
\begin{enumerate}
    \item another same-output deeper-origin relay that merely preserves the current reduced projection-response substrate structure \(\ReducedStructure\) and therefore merely reproduces the same full image core \(\CurrentBodyImageCore\) together with the same downstream benchmark-facing consequences;
    \item another restoration of the full substrate-body presentation as if the already-screened larger body freedom were again independently live;
    \item another replay of the older projection-operator core, the full displacement-response substrate law, or the larger pullback-body-operator, pullback-operator, pullback-quadratic, hidden-response, residual-normal-form, or pre-proto-section presentations as if those already-spent larger presentations were again the honest current burden;
    \item a manuscript that introduces a second operative metric, enlarges the ordinary matter law, or uses stronger promotion language.
\end{enumerate}
\end{proposition}

\begin{proof}
Item (1) fails the preceding criterion because it leaves the renewed live burden unchanged. Item (2) likewise fails because it restores a larger presentation whose benchmark-facing freedom is already screened out. Item (3) fails because it re-expands to larger already-spent presentations rather than deriving the reduced structure now known to force the full image core or proving a sharper theorem that forces that derivation. Item (4) violates the fixed claim boundary of the benchmark package and therefore cannot count as a disciplined next step on the active line.
\end{proof}

\begin{remark}
This screening statement is not a denial that those deeper or alternative manuscripts may matter strategically. It is a routing statement: they are not the default way to spend the next benchmark-facing move below the current frontier unless they do the same benchmark-facing work named above.
\end{remark}

\section{Conclusion}

The positive result of this note is routing discipline below the canonical-current-body-image-core necessity / uniqueness frontier. The present theorem remains the live benchmark-facing gain because it already proves that the full canonical current-body image core \(\CurrentBodyImageCore\) is canonically derived and uniquely forced once the reduced projection-response substrate structure \(\ReducedStructure\) is fixed. The earlier body-, relation-, response-, and shell-level theorems remain preserved main-line history because they removed those constituents one by one before the current cleaner full-image-core compression.

The scope remains narrow and honest. The benchmark package is unchanged: the one operative metric is still exactly \(\CompletedMetric\), the ordinary matter route is unchanged, there is no second operative metric, and the low-energy matter law is not enlarged. Nothing here upgrades adoption to derivation or licenses stronger promotion language.

The next open burden is therefore clear. The next benchmark-facing manuscript should derive or justify the reduced projection-response substrate structure \(\ReducedStructure\) from still more primitive explicit substrate structure, or prove a still sharper theorem that makes such a derivation unavoidable. Another same-output reduced-structure relay, another restoration of the full substrate-body presentation, or another replay of the older projection-operator core or the full displacement-response substrate law is not the clean next manuscript target at current scope.

\end{document}
